%%==========================================================================
\section{The Maintenance-Extended Model and Cumulative-Flight-Time Correction}
\label{sec:maintenance}
%%==========================================================================

\subsection{Maintenance problem definition}

Following \citet{khaled2018compact}, we consider type-A checks triggered after
at most $T^A_{\max}$ cumulative flight hours.  Checks take place at night
after the last flight on a given day in an airport with maintenance capacity.

Additional notation is collected in \Cref{tab:maint_notation}.

\begin{table}[t]
\centering\small
\caption{Additional notation for the maintenance model.}
\label{tab:maint_notation}
\begin{tabular}{ll}
\toprule
Symbol & Meaning \\
\midrule
$\D = \{1,\ldots,n_d\}$      & Planning days \\
$d_i$                         & Departure day of flight $i$ \\
$\bar d_i$                    & Arrival day of flight $i$ \\
$\MA \subseteq \A$            & Maintenance-capable airports \\
$\text{mcap}_{m,d}$           & Maintenance capacity of airport $m$ on day $d$ \\
$mc_{ijd}$                    & Cost of night maintenance at arrival airport of
                                flight $i$, aircraft $j$, day $d$ \\
$F_d^{\downarrow}(m)$         & Flights arriving at airport $m$ on day $d$ \\
$F(d,i)$                      & $\{i' : d_{i'}=d,\; \tD{i'} > \tI{i}\}$ \\
$T^A_{\max}$                  & Max cumulative flight minutes before type-A check \\
$d_{\max}$                    & Max calendar-day interval between checks \\
$M_{dd'}$                     & Big-$M$ for day pair $(d,d')$ \\
$F_{d,d'}$                    & Flights with departure day in $(d,d']$ \\
\bottomrule
\end{tabular}
\end{table}

\subsection{Additional decision variables}
\paragraph{C6, C7 — Boolean variables for maintenance.}
\begin{align*}
  z_{ijd} &= \begin{cases}1 & \text{night maintenance at arrival airport of } i,\;
               \text{aircraft } j,\;\text{day } d, \\ 0 & \text{else;}\end{cases}\\
  y_{jd}  &= \begin{cases}1 & \text{maintenance of aircraft } j \text{ on day } d,
               \\ 0 & \text{else,}\end{cases}
\end{align*}

The variable $z_{ijd}$ is instantiated only when flight $i$ arrives at a
maintenance-capable airport on day $d$. Thus its admissible domain is
$i\in\bigcup_{m\in\MA}F_d^{\downarrow}(m)$. We assume in this type-A model
that a night check fits within the overnight ground period.

\subsection{Extended model}

The objective becomes:
\begin{equation}\label{eq:obj_ext}
  \min\; \sum_{i \in \F}\sum_{j \in \Pc} c_{ij} x_{ij}
  + \sum_{d \in \D}\sum_{m \in \MA}\sum_{i \in F_d^{\downarrow}(m)}
    \sum_{j \in \Pc}
    mc_{ijd}\, z_{ijd}.
\end{equation}

Constraints~\eqref{eq:c0}--\eqref{eq:c5} are retained.  Additional
constraints:

\paragraph{C8 — Maintenance blocks same-day subsequent flights.}
\begin{equation}\label{eq:c8}
  z_{ijd} + x_{i'j} \le 1,
  \quad \forall d\in\D,\;m\in\MA,\;i\in F_d^{\downarrow}(m),\;
  i'\in F(d,i),\;j\in\Pc.
\end{equation}

\paragraph{C9 — Maintenance requires assignment.}
\begin{equation}\label{eq:c9}
  x_{ij} \ge z_{ijd},
  \quad \forall d\in\D,\;m\in\MA,\;i\in F_d^{\downarrow}(m),\;j\in\Pc.
\end{equation}

\paragraph{C10 — Capacity.}
\begin{equation}\label{eq:c10}
  \sum_{i \in F_d^{\downarrow}(m)} \sum_{j \in \Pc} z_{ijd}
  \le \text{mcap}_{m,d},
  \qquad \forall d \in \D,\; m \in \MA.
\end{equation}

\paragraph{C11 — Aggregation ($z \to y$).}
\begin{equation}\label{eq:c11}
  \sum_{m\in\MA}\sum_{i \in F_d^{\downarrow}(m)} z_{ijd} = y_{jd},
  \qquad \forall d \in \D,\; j \in \Pc.
\end{equation}

\paragraph{C12 — Maximum-spacing constraint.}
\begin{equation}\label{eq:c12}
  \sum_{r \in [d, d']} y_{jr} \ge 1,
  \qquad \forall d, d' \in \D \text{ s.t.\ } d' - d = d_{\max} - 1,\;
  \forall j \in \Pc.
\end{equation}

For the type-A-only model studied in this paper, C12 is not an additional
calendar-spacing requirement: type-A maintenance is governed by the exact
flight-hour state introduced in the cumulative-state formulation below. The displayed row is retained as the generic
day-indexed maintenance notation and becomes active only for a calendar-time
maintenance type, which is outside the scope of this paper.

\paragraph{C13 — At most one maintenance event per aircraft per day.}
\begin{equation}\label{eq:c13}
  \sum_{m\in\MA}\sum_{i \in F_d^{\downarrow}(m)} z_{ijd} \le 1,
  \qquad \forall j \in \Pc,\; d \in \D.
\end{equation}

For binary $y$, C13 is implied by C11; it is retained because it makes the
one-event-per-day rule explicit and remains useful in implementations that
relax or disaggregate the aggregation equality.

The original paper's cumulative-flight-time constraint is:
\paragraph{C14 — Paper-based cumulative-flight-time constraint.}
\begin{equation}\label{eq:c14_paper}
  \sum_{i \in F_{d,d'}} x_{ij}\,\tilde{t}_i
  \;\le\;
  T^A_{\max}
  + M_{dd'}\Bigl(2 - y_{jd} - y_{jd'}\Bigr)
  + M_{dd'}\!\sum_{r \in [d+1,d'-1]} y_{jr},
\end{equation}
for all $j \in \Pc$, all pairs $(d, d')$ with $2 \le d'-d \le d_{\max}-1$,
$d < n_d$, where $\tilde{t}_i$ is the flight duration of leg $i$.


\subsection{Impact of corrected C0--C5 feasibility on C8--C14}
\label{subsec:c2c4_impact_maintenance}

The correction introduced in \Cref{sec:basic} (non-overlap
constraint~\eqref{eq:c4_overlap}) changes the admissible assignment set for
$x$, and therefore the interpretation of all maintenance constraints layered on
top of $x$.

\paragraph{C8.}
No algebraic change is required in~\eqref{eq:c8}; however, with
\eqref{eq:c4_overlap} active, the blocking relation now operates on physically
realizable trajectories only. This removes spurious cases where maintenance
placement was evaluated on infeasible simultaneous departures.

\paragraph{C9--C12.}
Constraints~\eqref{eq:c9}--\eqref{eq:c12} remain structurally unchanged. Their
feasible region is tightened indirectly because $z$ and $y$ are linked to a
strictly smaller, physically valid $x$-space.

\paragraph{C13 and cumulative-hour rows.}
Hour-accumulation constraints are evaluated over a physically executable
$x$-space once~\eqref{eq:c4_overlap} is imposed. This correction is necessary,
but it does not by itself establish that a particular cumulative-hour algebraic form
correctly propagates the maintenance counter; that property is analysed below.

In summary, the C0--C5 correction primarily changes the domain over which the
maintenance constraints operate. It is independent of the separate algebraic
correction required for cumulative-hour propagation.

\subsection{Issues of day-indexed maintenance constraints C8--C13}

Usage of day-indexed maintenance variables $y_{jd}$ and $z_{ijd}$ is
operationally intuitive, gives clear modeling semantics, but it has the following drawbacks:
\begin{itemize}
\item The day-indexed formulation requires a large number of big-$M$ constants
  $M_{dd'}$ to cover all day pairs $(d,d')$ with $2\le d'-d\le d_{\max}-1$.
  This is a combinatorial number of rows, and the big-$M$ values are
  necessarily loose because they must cover all possible flight-hour
  accumulations over the horizon.

  \item The day-indexed formulation does not propagate the flight-hour state
  exactly.  The big-$M$ relaxation allows the flight-hour counter to be reset
  prematurely, leading to potential violations of the intended maintenance
  intervals. For example, if a maintenance night is scheduled on day $d$ and the next maintenance night is 
  scheduled on day $d'$, the big-$M$ constraints only enforce that the cumulative flight hours between $d$ and $d'$ do not exceed $T^A_{\max}$, but they do not prevent excessive flight hours before day $d$.

  \item The day-indexed formulation forces a specific maintenance schedule on each day for a given time zone, thus allowing to have two different maintenance checks within 24 hrs, but not in two different days within the same 24-hr period.

  \item The day-indexed formulation does not allow for a natural integration with the ordered flight formulation, which is more efficient and compact. The ordered flight formulation allows for a more direct representation of the maintenance events in relation to the flight sequence, leading to a more efficient model.

  \item The original cumulative-flight-time row (paper equation) in the day-indexed formulation is flawed, as it does not correctly enforce the flight-hour limit between two consecutive maintenance nights. This flaw can lead to infeasible solutions that violate the intended maintenance intervals.
\end{itemize}


\subsection{Legacy cumulative-flight-time row: original formulation and flaw}
\label{subsec:c11_flaw}

The original paper's cumulative-flight-time constraint \eqref{eq:c13_paper} is intended to enforce that the cumulative flight hours between two consecutive maintenance nights do not exceed $T^A_{\max}$. However, as shown in \Cref{lem:c11_flaw}, this constraint can fail to enforce the intended maintenance intervals in certain cases, particularly when there is no maintenance scheduled at the start of the interval.

\begin{lemma}[Flaw in~\eqref{eq:c13_paper}]
  \label{lem:c11_flaw}
  Consider two consecutive maintenance events at days $d$ and $d'$ with no
  check between them ($y_{jr} = 0$ for $r \in (d, d')$). Then the RHS
  of~\eqref{eq:c13_paper} equals $T^A_{\max} + M_{dd'}(2 - 1 - 1) +
  M_{dd'} \cdot 0 = T^A_{\max}$, which correctly imposes the flight-hour
  limit between the two checks.

  However, if $y_{jd} = 0$ (no maintenance at the start day) and $y_{jd'} =
  1$ (maintenance only at the end day), the RHS becomes $T^A_{\max} +
  M_{dd'}(2 - 0 - 1) = T^A_{\max} + M_{dd'}$, which is vacuously large and
  imposes \emph{no limit} on cumulative flight hours before the first check of
  the period. Hence the constraint fails to enforce the check-interval bound
  for the opening sub-period.
\end{lemma}

\begin{proof}
  Direct substitution of $y_{jd} = 0$, $y_{jd'} = 1$, $\sum_{r} y_{jr} = 0$
  into~\eqref{eq:c13_paper} gives the stated RHS. Since $M_{dd'}$ is chosen
  as the maximum possible cumulative flight time over the horizon (a large
  positive constant), the resulting constraint is trivially satisfied regardless
  of the actual flight hours flown.
\end{proof}

\subsection{Legacy cumulative-flight-time row: endpoint-split strengthening and limitation}
\label{subsec:c11_corrected}

An endpoint-split strengthening considered for replacing
\eqref{eq:c13_paper} consists of two constraints, each relaxed by one endpoint:

\begin{equation}\label{eq:c11_split_a}
  \sum_{i \in F_{d,d'}} x_{ij}\,\tilde{t}_i
  \;\le\;
  T^A_{\max}
  + M_{dd'}\, y_{jd}
  + M_{dd'}\!\sum_{r \in [d+1,d'-1]} y_{jr},
\end{equation}
\begin{equation}\label{eq:c11_split_b}
  \sum_{i \in F_{d,d'}} x_{ij}\,\tilde{t}_i
  \;\le\;
  T^A_{\max}
  + M_{dd'}\, y_{jd'}
  + M_{dd'}\!\sum_{r \in [d+1,d'-1]} y_{jr},
\end{equation}
for all $j \in \Pc$, $(d,d') \in \D^2$ with $2 \le d'-d \le d_{\max}-1$,
$d < n_d$.

\begin{lemma}[The endpoint split is not sufficient]
  \label{lem:c11_split_limit}
  With $F_{d,d'}$ defined as the flights departing in $(d,d']$,
  constraints~\eqref{eq:c11_split_a}--\eqref{eq:c11_split_b} do not necessarily
  bound all flying time between two consecutive maintenance nights.
\end{lemma}

\begin{proof}
  Let $T^A_{\max}=10$, $M_{dd'}=20$, $d_{\max}\geq4$, and consider days
  $1,\ldots,4$ with
  $y_{j1}=y_{j4}=1$ and $y_{j2}=y_{j3}=0$. Assign 10 units of flying time on
  day 2 and 10 units on day 4. For $(d,d')=(1,3)$, the accumulated time is 10
  and the second split inequality is active. For $(2,4)$, the accumulated time
  is 10 and the first split inequality is active. For $(1,4)$, the accumulated
  time is 20, but both inequalities have right-hand side
  $T^A_{\max}+M_{14}=30$ because both endpoints contain checks. Thus every
  split inequality is satisfied although 20 units are flown between the
  consecutive checks, exceeding the limit 10.
\end{proof}

The original and split formulations are therefore not ordered by strength:
the original row is active when both endpoints contain checks, whereas the
split rows are active when at least one corresponding endpoint indicator is
zero. Neither pattern alone gives an exact cumulative counter.

\subsection{Exact cumulative-state formulation for cumulative flight hours}
\label{subsec:c11_state}

Define the flying time assigned to aircraft $j$ on day $d$ by
\paragraph{C14a Daily flight time}
\begin{equation}
  v_{jd}=\sum_{i\in\F:\,d_i=d}\tilde{t}_i x_{ij}.
  \label{eq:daily_flight_time}
\end{equation}
Let $h_{jd}\geq0$ be the accumulated type-A flying time immediately after the
flights of day $d$ and before maintenance that night. Set
$h_{j0}=\Phi^A_j$ and $y_{j0}=0$. For consecutive days, impose
\paragraph{C14b-C14f Accumulated flight time constraints}
\begin{align}
  h_{jd}
  &\geq h_{j,d-1}+v_{jd}-M^H y_{j,d-1},
  \label{eq:c11_state_lower}\\
  h_{jd}
  &\leq h_{j,d-1}+v_{jd}+M^H y_{j,d-1},
  \label{eq:c11_state_upper}\\
  h_{jd}
  &\geq v_{jd},
  \label{eq:c11_state_reset_lower}\\
  h_{jd}
  &\leq v_{jd}+M^H(1-y_{j,d-1}),
  \label{eq:c11_state_reset_upper}\\
  0\leq h_{jd}&\leq T^A_{\max},
  \qquad \forall j\in\Pc,\ d\in\D,
  \label{eq:c11_state_bound}
\end{align}
where $M^H\geq T^A_{\max}$ is valid because
$h_{j,d-1}\leq T^A_{\max}$. If $y_{j,d-1}=0$, the first two inequalities give
$h_{jd}=h_{j,d-1}+v_{jd}$. If $y_{j,d-1}=1$, the reset inequalities give
$h_{jd}=v_{jd}$. The upper bound then enforces the flight-hour threshold on
every day, including the flights operated immediately before a night check.

\subsection{Pre-horizon accumulated hours}
\label{subsec:c11_pre_horizon}

The original paper does not account for flying time accumulated before the
planning horizon. In the cumulative-state formulation this history is included
exactly through the initial condition
\begin{equation}\label{eq:c11_pre_horizon}
  h_{j0}=\Phi^A_j,
  \qquad 0\leq\Phi^A_j\leq T^A_{\max},\quad \forall j\in\Pc.
\end{equation}
The condition is required for every aircraft, including $\Phi^A_j=0$; no
separate opening-window big-$M$ inequality is needed.

\subsection{Integrality constraints}
\begin{equation}
  x_{ij} \in \{0,1\},\quad
  z_{ijd} \in \{0,1\},\quad
  y_{jd} \in \{0,1\},
  \qquad
  \begin{aligned}
    &(i,j)\in\F\times\Pc,\\
    &d\in\D,\quad
      i\in\bigcup_{m\in\MA}F_d^{\downarrow}(m)\ \text{for }z_{ijd}.
  \end{aligned}
\end{equation}

The variables $h_{jd}$ are continuous and satisfy
$0\leq h_{jd}\leq T^A_{\max}$ as imposed in~\eqref{eq:c11_state_bound}.

As we see, the fixes to the original paper's maintenance model are twofold: (i) the C0--C5 correction, which ensures that the flight assignment is physically feasible, and (ii) the cumulative-flight-time row correction.
Also, the cumulative-flight-time correction is implemented in two ways: (i) the endpoint-split strengthening, which is a partial fix, and (ii) the exact cumulative-state formulation, which is a complete fix. The exact cumulative-state formulation correctly propagates the flight-hour state and enforces the maintenance intervals as intended.
The complexity of the corrected model is higher than the original model, because we need to track the cumulative flight hours explicitly that require more $O(n_f D)$ variables, 
but it is necessary to ensure that the solutions are physically feasible and respect the maintenance requirements.
day-indexed model is higher than the original
model, because the cumulative flight-hour state $h_{jd}$ adds
$O(n_p n_d)$ continuous variables,

The additional complexity is justified by the need to accurately model the maintenance requirements and ensure that the resulting flight schedules are feasible in practice.

\subsection{Rigorous feasibility theorem for C1--C14}
\label{subsec:feasibility_theorem_c1_c14}

We now state the corrected sufficiency result for the type-A maintenance model.
Consider the constraint family
\eqref{eq:c0}--\eqref{eq:c5}, \eqref{eq:c8}--\eqref{eq:c11}, \eqref{eq:c13},
\eqref{eq:c11_state_lower}--\eqref{eq:c11_state_bound}, and the stated
variable domains. Constraint C12 is inactive for the type-A-only model:
hour spacing is enforced by the cumulative-state rows rather than by
calendar-day covering.

\begin{theorem}[Physical feasibility under corrected C1--C14]
  \label{thm:physical_feasibility_c1_c14}
  Any solution satisfying the above constraints, with $x$, $z$, and $y$
  binary and $h$ continuous, induces for each aircraft $j$ a physically
  executable trajectory over the planning horizon:
  flights assigned to $j$ are temporally non-overlapping, satisfy airport
  continuity and turn time, and admit a day-by-day maintenance schedule
  satisfying capacity and hour-spacing rules.
\end{theorem}

\begin{proof}
  Fix an aircraft $j$ and let
  $S_j=\{i\in\F: x_{ij}=1\}$ ordered by nondecreasing departure time.

  \textbf{(i) Flight assignment and continuity.}
  By~\eqref{eq:c1}, every flight is assigned to exactly one aircraft.
  By~\eqref{eq:c2}--\eqref{eq:c3}, each departure assigned to $j$ has sufficient
  prior balance at its departure airport, with the initial unit at $a^0_j$.
  Therefore airport continuity holds along $S_j$.

  \textbf{(ii) Temporal executability.}
  For any overlapping pair $\{i_1,i_2\}\in\mathcal{O}$,
  \eqref{eq:c4_overlap} enforces $x_{i_1j}+x_{i_2j}\le 1$.
  Hence no two flights assigned to $j$ overlap once turn time is included.
  Combined with (i), this yields a realizable single-aircraft flight timeline.

  \textbf{(iii) Maintenance event consistency.}
  By~\eqref{eq:c9}, maintenance trigger variables can only be activated on
  assigned flights. By the aggregation equality~\eqref{eq:c11}, each day's
  maintenance trigger is represented by the single aggregate indicator $y_{jd}$;
  by the daily cap~\eqref{eq:c13}, at most one maintenance event is attached to
  aircraft $j$ on day $d$.

  \textbf{(iv) Maintenance capacity and blocking.}
  By~\eqref{eq:c10}, airport-day capacities are respected. By~\eqref{eq:c8}, if
  maintenance is triggered after flight $i$ on day $d$, then any subsequent
  same-day flight is forbidden for the same aircraft. Under the stated
  overnight-duration assumption, the maintenance action is operationally
  executable.

  \textbf{(v) Hour/day maintenance admissibility.}
  By~\eqref{eq:c11_state_lower}--\eqref{eq:c11_state_bound}, the cumulative
  flight-time state increases by the flying time assigned each day, resets
  after a maintenance night, and never exceeds $T^A_{\max}$.

  Steps (i)--(v) show that every aircraft has a consistent, non-overlapping
  flight sequence with feasible maintenance placement and resource usage.
  Therefore the global solution is physically feasible.
\end{proof}
